\section{Current research priority and decision criteria}
\label{sec:outlook}

The finite WZW pilot is complete: it has a boundary-selected initial block,
a fully normalized observable, closed deterministic evolution, explicit
fusion and composition, and a spatial-smearing test. Its constant-driver
balance is in a specified tensor pairing. Its proposed extension is a slit
sewing map on the boundary affine modules, retaining descendants and
endpoint fields, with a physical reflected pairing. Such a construction
must reproduce the primary matrix elements above and explain its relation
to both metrics in Section~\ref{sec:pairings}.

For arithmetic matching, the direct source calculation now gives a sharper
starting point than further sewing without an arithmetic specification.
The next bounded task is a cumulative boundary or canonical response for
$K_\omega$ or $Z_\omega$, starting from the complete source on a safe
line. Its deliverable is an operator dictionary in a fixed Hilbert space
and an explicit norm balance, including the cost of removing the weight
in \eqref{eq:weighted_bound}. Three tests make that task concrete:
\begin{enumerate}
\item Recover the identity initial operator and the full local, gamma and
rational terms on $0<L<\log2$, before any prime delay is visible.
\item On crossing $L=\log2$, recover the exact first-delay coefficient
$-\sqrt2\log2\cosh(\omega\log2)$ in the logarithmic generator,
with control of every boundary or storage term.
\item Produce a cumulative positive balance compatible with
\eqref{eq:increment_counterexample}, rather than requiring positive
instantaneous multiplicative increments.
\end{enumerate}

Canonical-system work already provides a related comparison. Suzuki
\cite{Suzuki} studies
$\Theta_\omega(z)=K_\omega(-\ii z)$, its meromorphic innerness and
the associated canonical systems. His explicit unconditional construction
has $\omega>1$. This is different from unconditional innerness for
$\omega\ge1/2$, and from the safe source translation $\eta\ge1/2$.
Innerness for every positive arithmetic shift is RH-equivalent. Canonical
depth must therefore not be silently identified with $\omega$ or $\tau$.

The modular-surface scattering realization at $\omega=1/2$, already
recorded in the parent investigation, remains an endpoint comparison. It
does not supply a continuous family at smaller shifts. Likewise,
postulating the measure \eqref{eq:atomic_measure} without RH would assume
the desired positive-real property instead of deriving it.

An independent realization must specify its operator and positive norm
without assuming the contractivity it is meant to explain. The present
work provides an explicit physical benchmark, an exact arithmetic source
for comparison, and several rejection tests for proposed identifications.
It does not establish a matching physical transfer, its unweighted
cumulative norm identity, or an all-depth continuation.
